% Learning InitechBase, Chapter 6: Calculating (AN-16).  Every panel
% replayed (examples/l6-*.dot); sources cmd/flow.c (memvars), core/eval.c,
% core/fn_builtins.c, cmd/set.c (DATE, CENTURY, DECIMALS).
\chapter{Calculating}
\Gates{test-pubs-examples, test-xbase-eval, test-xbase-fn-a, test-xbase-fn-b, test-interp-setfmt}
InitechBase calculates with the values of fields and with values of its own:
memory variables.  In this lesson you store values in memory variables,
calculate with numbers, character strings and dates, and use the functions
that convert one kind of value into another.

\begin{ChapterCovers}
\item Memory variables
\item Arithmetic
\item Functions of character strings
\item Converting numbers to characters
\item Dates and date arithmetic
\item SET DATE, SET CENTURY and SET DECIMALS
\item Choosing a value with \fn{IIF()}
\end{ChapterCovers}

\ATEHead{Preparing for This Chapter}
Start InitechBase with a fresh copy of the sample files.

\ATEHead{Memory Variables}
A memory variable is a named value that InitechBase keeps for you until you
leave it.  \cmd{STORE} puts a value in one or more memory variables; the
equals sign does the same for one.
\begin{Steps}
\item Type \type{STORE 0 TO MCOUNT}.
\item Type \type{MNAME = 'INITECH SYSTEMS'}, then \type{MRAISE = 4.5}.
\item Type \type{? MCOUNT, MNAME, MRAISE}.
\end{Steps}
\DotExample{examples/l6-memvar}
The names begin with M so that they cannot be taken for the names of fields.

\ATEHead{Arithmetic}
The arithmetic operators are \texttt{+}, \texttt{-}, \texttt{*} and \texttt{/};
parentheses group a calculation.  Here is Peter Gibbons' salary after a raise
of 4.5 per cent, and the amount of the raise:
\begin{Steps}[resume]
\item Type \type{USE STAFF}, then \type{GO 2}.
\item Type \type{MRAISE = 4.5}.
\item Type \type{MNEW = SALARY * (1 + MRAISE / 100)}.
\item Type \type{? MNEW, MNEW - SALARY}.
\end{Steps}
\DotExample{examples/l6-arith}

\ATEHead{Character Strings}
Functions work on values and return a value.  \fn{UPPER()} and \fn{LOWER()}
change the case of letters, \fn{LEN()} counts characters, and \fn{TRIM()}
removes trailing blanks:
\DotExample{examples/l6-strings}
\fn{SUBSTR()} takes part of a string, and \fn{LEFT()} its first characters:
\DotExample{examples/l6-substr}

\ATEHead{Numbers and Characters}
InitechBase does not join a number to a character string.  \fn{STR()} turns the
number into characters first: its arguments are the number, the width and the
decimal places.  \fn{ROUND()} rounds a number and \fn{INT()} drops its fraction.
\DotExample{examples/l6-numbers}

\ATEHead{Dates}
\fn{DATE()} is today's date.  Subtracting one date from another gives the
number of days between them; adding a number of days to a date gives a date.
\fn{YEAR()}, \fn{CMONTH()} and \fn{CDOW()} give the year, the name of the month
and the name of the day.
\begin{Steps}[resume]
\item Type \type{USE STAFF}, then \type{GO 9}.
\item Type \type{? DATE(), HIRED, DATE() - HIRED}.
\item Type \type{? YEAR(HIRED), CMONTH(HIRED), CDOW(HIRED)}.
\end{Steps}
\DotExample{examples/l6-dates}
The complete years of service of Andrew Delgado:
\DotExample{examples/l6-years}

\ATEHead{The Form of Dates}
SET DATE changes the order of the day, month and year; SET CENTURY ON shows the
year with four digits.
\DotExample{examples/l6-setdate}
Two digits are the standard: a date typed with a two-digit year is a date of
the nineteen hundreds.

\ATEHead{Decimal Places}
The result of a division is shown with two decimal places unless SET DECIMALS
says otherwise:
\DotExample{examples/l6-decimals}

\ATEHead{Choosing a Value}
\fn{IIF()} returns its second argument when the first is true and its third
when it is false:
\DotExample{examples/l6-iif}

\ATEHead{Where Do You Go From Here?}
You can now calculate with fields, memory variables and functions.  Chapter~6
of \doctitle{Using InitechBase} describes every function.  In Chapter~7 you
print a report.

\ATEHead{Summary}
\begin{enumerate}
\item \cmd{STORE} \xp{value} \kw{TO} \xp{names}, or \xp{name} \texttt{=}
  \xp{value}, stores values in memory variables.
\item \texttt{+ - * /} and parentheses calculate; \fn{ROUND()} and \fn{INT()}
  shorten a number.
\item \fn{UPPER()}, \fn{LOWER()}, \fn{LEN()}, \fn{TRIM()}, \fn{SUBSTR()} and
  \fn{LEFT()} work on characters; \fn{STR()} turns a number into characters.
\item \fn{DATE()} is today; dates can be subtracted, and days added to them.
\item SET DATE, SET CENTURY and SET DECIMALS change how values are shown.
\item \fn{IIF()} chooses between two values.
\end{enumerate}
